\section{A dividend is a date, not a rate}
\label{sec:deamdivs}

So far the tree has run on a smooth carry: the pair $(r,q_d)$ implied
by Chapter~6's fitted forward, with the dividend yield $q_d$ draining
value out of the spot continuously.  For the \emph{European} leg that
is often good enough --- a European price cares about the forward at
expiry, and Chapter~6 showed how to make any dividend convention
reproduce the right forward.  The \emph{stopping right} is pickier.
\Cref{sec:deamwhere} located the call-side premium in capturing a
dividend by exercising \emph{before its ex-date} --- the last moment
the stock still carries the payment.  A date, not a rate.  This
section teaches the tree about dates, and measures what smearing them
into a rate gets wrong.

The obvious construction fails, instructively.  Take the lattice of
\cref{sec:deamtree} and simply subtract the cash amount $d_1$ ---
Chapter~6's subscripted dividend symbol, not the Black argument
$d_\pm$ --- from every node at the ex-date.  Recombination dies on the spot: follow the
two two-step paths that should meet, with an ex-date between the
steps, writing $\mathrm{u}=e^{\sigma\sqrt{\Delta t}}$ just for this
display:
\[
\text{up, drop, then down: } (S\mathrm{u}-d_1)/\mathrm{u} = S -
d_1/\mathrm{u},
\qquad
\text{down, drop, then up: } (S/\mathrm{u}-d_1)\,\mathrm{u} = S -
d_1\mathrm{u} .
\]
The multiplicative moves scale the subtracted cash differently on
each path, the two nodes no longer coincide, and the lattice explodes
from $n+1$ nodes per layer into $2^{N_t}$ distinct paths --- from a
few hundred values to more than there are atoms in a spreadsheet.

The classical repair is the \emph{escrowed spot}
\citep{Hull2018}.  Split the stock into the part that will be paid out
and the part that will not: set aside (``escrow'') the present value
of the scheduled dividends, and let only the remainder diffuse.  The
diffusing base starts at
\[
X_0 \;=\; S-\mathrm{PV}_0,
\qquad
\mathrm{PV}_0=\textstyle\sum_i d_i\,e^{-r t_i},
\]
with cash amounts $d_i$ at ex-date year fractions $t_i\le t$, exactly
Chapter~6's schedule objects.  The base lattice is purely
multiplicative --- it recombines, cash never touches it --- and the
\emph{true} spot at any node is reconstructed when needed: base value
plus the present value of the dividends \emph{not yet paid} at that
layer's date.  The exercise comparison in \eqref{eq:deamroll} tests
intrinsic against that true cum-dividend spot, so the holder's
decision sees every upcoming payment; at expiry every scheduled
dividend has been paid and the payoff reads the bare base.  One
point of precision: under escrow the thing that diffuses lognormally is the
dividend-stripped spot, so the model's spot distribution is a
\emph{displaced} lognormal --- a slightly different law than the
smooth-carry tree at the same forward.  That is a modelling choice
inside the premium, and it goes in the contract with the others.

One consistency condition ties the schedule to Chapter~6, and it must
hold \emph{exactly}: the escrowed tree's forward,
$(S-\mathrm{PV}_0)\,e^{rt}$, has to reproduce the fitted forward $F$
--- otherwise the premium machine and the parity line disagree about
the same chain.  When the forecast schedule misses by a little, the
repair follows Chapter~6's escrow lesson --- the \emph{timing} is the
forecast's, the \emph{amounts} bend to the market --- by rescaling
every cash amount by the common factor
\begin{equation}
\frac{S-F e^{-rt}}{\sum_i d_i\,e^{-r t_i}}\,,
\label{eq:deamalpha}
\end{equation}
which is exactly the ratio that makes the escrowed forward land on
$F$.  (A nonsensical factor --- negative, or far from one --- is a
sign the schedule itself is wrong, and rejects it.)

Now the measurement: what does the date-aware tree know that the
smeared one cannot?  \Cref{fig:deamescrow} prices a deliberately
heavy payer --- one \$6 dividend on a \$100 stock, ex-date at $0.35$
years, against a $2\%$ rate --- both ways.  In panel~(a), fix the
expiry at $t=0.5$ and hand both models the \emph{same} forward
(\$\MacDeamEscrowFwd; the matching smooth yield is
$q_d=\MacDeamEscrowQeqPct\%$).  The premiums still disagree ---
everywhere, and by up to \MacDeamEscrowRatioMax$\times$ where both
are meaningfully positive.  Matching the forward is not enough,
because the premium prices the \emph{timing} of the payout: a lump
of value leaving at one date is worth more to an exerciser --- who
can position just ahead of it --- than the same value bleeding out
smoothly.

Panel~(b) is the sharper exhibit.  Fix the in-the-money $K=85$ call
and sweep the \emph{expiry} across the ex-date.  The date-aware
premium (blue) is \emph{exactly zero} --- machine zero,
$\MacDeamEscrowCashPreExMax$ --- for every expiry before $0.35$:
an option that expires before the ex-date has no dividend inside its
life to capture, so \cref{thm:deammerton} applies verbatim.  At the
ex-date the premium does not fade in; it \emph{jumps}.  The reason is
worth one plain sentence: for an expiry just past the date, the
holder can exercise the instant before the stock goes ex and collect
the full cum-dividend intrinsic, while the European twin must ride
through the drop --- so crossing the date hands the best plan, at
once, roughly the dividend times the probability of finishing in the
money (the panel's jump is a little under \$5 of the \$6, this call
being deep but not certain).  The smeared comparator (green, the schedule's
one-year-equivalent flat yield $q_d=\MacDeamEscrowFlatQPct\%$)
misses both regimes: before the date it \emph{invents}
\$\MacDeamEscrowFlatPreExMax\ of premium where the truth is exactly
zero, and after the date it underprices the lump it spread out.  A
dividend forecast, for the stopping right, is a calendar --- the
amounts matter, and the dates matter separately.

\begin{figure}[htbp]
\centering
\paperfig{\textwidth}{fig_deam_escrow.pdf}
\caption{One \$6 cash dividend, ex-date $0.35$ years, $r=2\%$,
$\sigma=25\%$.  (a)~Calls at $t=0.5$ under the escrowed schedule
versus the smooth yield matched to the \emph{same} forward
(\$\MacDeamEscrowFwd): equal forwards, unequal premiums --- up to
\MacDeamEscrowRatioMax$\times$ --- because the premium prices the
payout's timing.  (b)~The $K=85$ call across expiries: the date-aware
premium is exactly zero before the ex-date (nothing inside the
option's life to capture: \cref{thm:deammerton} verbatim) and jumps
across it; the flat smear ($q_d=\MacDeamEscrowFlatQPct\%$) invents
\$\MacDeamEscrowFlatPreExMax\ of premium where the truth is zero and
underprices the lump after.}
\label{fig:deamescrow}
\end{figure}

For the frozen SPY chains that thread this book, the chapter's own
real-data figures run the smooth-carry tree with the
forward-consistent $q_d$: an ETF's quarterly dividends are small, and
the fitted forward already carries their information.  The residual
cost of that convention is measured, not assumed, in the next section
--- a few basis points of forward.  For a single name with heavy
dates, the escrowed tree is the correct default, and either choice
rides into $\sigma^*$ as part of the model.

We now have the machine, the subtraction, and the carry.  One debt to
Chapter~6 remains --- the forward that all of this runs on was itself
fitted from American quotes.
